\documentclass[11pt,a4paper]{article}
\usepackage[margin=2.5cm]{geometry}
\usepackage{times}
\usepackage{latexsym}
\usepackage{booktabs}
\usepackage{amsmath}
\usepackage{amssymb}
\usepackage{graphicx}
\usepackage{multirow}
\usepackage{xcolor}
\usepackage{enumitem}
\usepackage{hyperref}

\newcommand{\pp}{\,pp}
\newcommand{\cmark}{\checkmark}

\title{Layer Selection Matters: How Metric Choice Determines\\Hidden-State Probe Effectiveness for Reasoning Model Safety Monitoring}

\author{Anonymous}

\date{}

\begin{document}
\maketitle

\begin{abstract}
Reasoning-capable language models can ``self-jailbreak'' during extended chain-of-thought, gradually abandoning safety reasoning until a \emph{commitment point} after which harmful content is generated.
We present the first systematic evaluation of detection methods for this safety commitment across four models spanning two architectures (Llama, Qwen), two training methods (R1-distillation, STILL), and two scales (7--8B, 32B).
Our primary methodological finding is that the layer selection metric \texttt{bal\_acc $\times$ crossing\_rate} systematically biases toward shallow layers at 32B scale, entirely masking hidden-state (HS) probe advantages---flipping results from 5/12 to \textbf{12/12} significant comparisons under Holm--Bonferroni correction when corrected layers are used.
At correctly selected layers, HS probes significantly outperform text-based probes on all four models (precision $+$15--24$\pp$, $p < 0.01$), with temporal context providing universal improvement.
Six alternative detection approaches---Mahalanobis distance, RepE direction probing, CCS, MLP probes, BOCPD, and CUSUM---fail comprehensively across all models, establishing supervised temporal probing with principled layer selection as the most effective approach for monitoring reasoning-model safety commitment.
\end{abstract}

%----------------------------------------------------------------------
\section{Introduction}
\label{sec:intro}
%----------------------------------------------------------------------

Reasoning-capable large language models (LLMs) such as DeepSeek-R1 and QwQ exhibit a novel safety failure mode: during extended chain-of-thought (CoT), the model's own reasoning can gradually erode safety constraints, leading to a ``self-jailbreak'' where safety reasoning permanently ceases at a \emph{commitment point} and harmful content follows \cite{huang2025selfJailbreaking,wen2025harmThoughts}.
Unlike adversarial input attacks on instruction-tuned models, the threat originates from the model's own reasoning process.

Existing hidden-state detectors such as TrajGuard \cite{zhang2025trajguard} and SafeSwitch achieve high accuracy against adversarial inputs but operate at token level and have never been evaluated on the distinct dynamics of reasoning-model self-jailbreak.
Text-level CoT monitors \cite{wang2025safeRemind,zheng2025safeThink} are vulnerable to pathological CoT where surface-level safety language persists while the model's internal commitment has already shifted.
No prior work systematically compares detection modalities---hidden states versus text---for reasoning-model safety commitment.

We address this gap with a comprehensive evaluation spanning four models: R1-8B (Llama-8B, R1-distilled), OpenThinker-7B (Qwen-7B, STILL-trained), QwQ-32B (Qwen-32B, native reasoning), and R1-32B (Qwen-32B, R1-distilled).
Our contributions are threefold:

\begin{enumerate}[nosep,leftmargin=*]
\item \textbf{Methodological:} We discover that the composite layer selection metric \texttt{bal\_acc $\times$ crossing\_rate} introduces a systematic bias toward shallow layers at 32B scale, completely masking the hidden-state advantage.
A sensitivity analysis across 5~metrics $\times$ 4~models reveals this as a researcher degree of freedom affecting any study using threshold-frequency--contaminated metrics (\S\ref{sec:layer_selection}).
\item \textbf{Empirical:} At correctly selected layers, HS probes significantly outperform text probes on all four models (12/12 Holm--Bonferroni, precision $+$15--24$\pp$), with temporal context universally improving both modalities (\S\ref{sec:hs_vs_text}).
\item \textbf{Negative results:} Six alternative detection methods fail comprehensively, establishing baselines and ruling out entire method families (\S\ref{sec:negative}).
\end{enumerate}

%----------------------------------------------------------------------
\section{Related Work}
\label{sec:related}
%----------------------------------------------------------------------

\paragraph{Reasoning model safety.}
Self-Jailbreaking \cite{huang2025selfJailbreaking} characterizes how reasoning models abandon safety constraints during CoT.
HarmThoughts \cite{wen2025harmThoughts} provides taxonomy of harmful reasoning patterns.
Neither proposes deployable detectors.

\paragraph{Hidden-state safety monitoring.}
TrajGuard \cite{zhang2025trajguard} uses Mahalanobis distance on hidden states for token-level adversarial detection; SIREN \cite{peng2025siren} monitors internal representations for jailbreak attempts.
Both target input-level attacks on instruction-tuned models, not reasoning-model self-jailbreak.

\paragraph{Representation probing for safety.}
Representation Engineering (RepE) \cite{zou2023repE} extracts safety-relevant directions via contrastive pairs.
CCS \cite{burns2023ccs} identifies truth directions unsupervised.
Linear probing of safety-relevant features has been studied extensively \cite{li2024probing}, but never for temporal commitment detection in reasoning models.

\paragraph{Change-point detection.}
BOCPD \cite{adams2007bocpd} and CUSUM detect distributional shifts in sequential data.
Their applicability to probe score sequences in reasoning traces has not been evaluated.

%----------------------------------------------------------------------
\section{Method}
\label{sec:method}
%----------------------------------------------------------------------

\subsection{Task: Safety Commitment Detection}

Given a reasoning trace $\mathbf{x} = (x_1, \ldots, x_T)$ partitioned into sentences, each sentence $x_t$ has a binary safety label $y_t \in \{0, 1\}$ (safe/unsafe).
The \emph{commitment point} (CP) is the last sentence after which all subsequent sentences are unsafe.
The detection task is to identify crossing from safe to committed-unsafe in real time.

\subsection{Probe Architectures}

\paragraph{Static linear probe.}
For each sentence $x_t$, we extract the hidden-state vector $\mathbf{h}_t^{(\ell)} \in \mathbb{R}^d$ at layer $\ell$ (last token pooling).
A logistic regression classifier is trained on $\mathbf{h}_t^{(\ell)}$ to predict $y_t$.

\paragraph{Temporal linear probe.}
We concatenate a window of $W$ consecutive hidden states: $\mathbf{z}_t = [\mathbf{h}_{t-W+1}^{(\ell)}; \ldots; \mathbf{h}_t^{(\ell)}] \in \mathbb{R}^{Wd}$, and train logistic regression on $\mathbf{z}_t$.
Temporal context captures confidence dynamics around the commitment point.

\paragraph{Text baseline.}
Sentence embeddings from all-MiniLM-L6-v2 replace hidden states.
The same static and temporal probe architectures are applied, enabling controlled modality comparison.

\paragraph{GRU sequence model.}
A GRU processes hidden-state sequences from multiple layers to capture temporal dynamics beyond fixed windows.

\subsection{Layer Selection Sensitivity Analysis}
\label{sec:layer_selection}

The choice of metric for selecting the optimal probe layer constitutes a critical researcher degree of freedom.
We evaluate five layer selection criteria across all four models:

\begin{enumerate}[nosep,leftmargin=*]
\item \texttt{bal\_acc $\times$ crossing\_rate}: composite metric penalizing low detection frequency;
\item \texttt{bal\_acc}: pure balanced accuracy (threshold-free);
\item \texttt{precision}: positive predictive value at the detection threshold;
\item \texttt{inverse FPR}: complement of pre-commitment-point false positive rate;
\item \texttt{composite}: additive combination of precision and inverse FPR.
\end{enumerate}

We discover that \texttt{bal\_acc $\times$ crossing\_rate} systematically selects shallow layers for 32B models because shallow layers have compressed representations (cosine distances 1.4--2.7$\times$ smaller), producing noise-driven threshold crossings that inflate the crossing rate.
Threshold-dependent metrics (precision, inverse FPR, composite) converge on deep layers; pure balanced accuracy selects intermediate layers.

\subsection{Statistical Framework}

All HS--text comparisons use null-centered two-sided bootstrap (10{,}000 resamples).
Family-wise error rate is controlled via Holm--Bonferroni step-down across 4~models $\times$ 3~metrics $=$ 12 tests at $\alpha = 0.05$.
Leave-one-out cross-validation assesses layer selection stability.

%----------------------------------------------------------------------
\section{Experimental Setup}
\label{sec:setup}
%----------------------------------------------------------------------

\paragraph{Models.}
We evaluate four reasoning models forming a 2$\times$2 architecture--scale matrix:
R1-8B (Llama-8B, R1-distilled, $n_{\text{test}}=43$),
OT-7B (Qwen-7B, STILL-trained, $n_{\text{test}}=47$, 228 total traces),
QwQ-32B (Qwen-32B, native, $n_{\text{test}}=19$--20),
R1-32B (Qwen-32B, R1-distilled, $n_{\text{test}}=20$).

\paragraph{Data.}
Reasoning traces are collected on safety-critical prompts (AdvBench-derived).
Each trace is sentence-segmented and annotated with binary safety commitment labels.
Train/val/test split: 60/20/20.

\paragraph{Metrics.}
We report balanced accuracy (BA), precision, pre-CP false positive rate (FPR), coverage (fraction of traces with at least one crossing), and lead time (steps before the commitment point at first detection).
All reported at sentence granularity.

%----------------------------------------------------------------------
\section{Results}
\label{sec:results}
%----------------------------------------------------------------------

\subsection{Layer Selection Metric Sensitivity}
\label{sec:layer_results}

Table~\ref{tab:layer_selection} presents the layer selected by each metric across all four models.
The biased metric \texttt{bal\_acc $\times$ crossing\_rate} selects extremely shallow layers for 32B models---L2 (3\% depth) for QwQ-32B and L10 (16\%) for R1-32B---while threshold-dependent metrics converge on deep layers: L60 (94\%) and L63 (98\%) respectively.

\begin{table}[t]
\centering
\small
\begin{tabular}{lcccc}
\toprule
\textbf{Metric} & \textbf{R1-8B} & \textbf{OT-7B} & \textbf{QwQ} & \textbf{R1-32B} \\
\midrule
BA $\times$ CR & L0 & L0 & L2 (3\%) & L10 (16\%) \\
Pure BA & L12 & L9 & L30 (47\%) & L47 (73\%) \\
Precision & L16 & L16 & L60 (94\%) & L63 (98\%) \\
Inv.\ FPR & L16 & L20 & L60 (94\%) & L63 (98\%) \\
Composite & L16 & L12 & L60 (94\%) & L63 (98\%) \\
\bottomrule
\end{tabular}
\caption{Layer selected by each metric per model. Percentages indicate depth within the 64-layer 32B models. The biased metric (BA $\times$ CR) selects shallow layers for 32B models; three threshold-dependent metrics converge on deep layers.}
\label{tab:layer_selection}
\end{table}

The impact of this bias on downstream conclusions is dramatic.
Table~\ref{tab:biased_vs_corrected} shows that at biased layers, the HS advantage over text is non-significant for both 32B models ($p = 0.050$--$0.346$); at corrected layers, significance flips to $p < 0.001$.

\begin{table}[t]
\centering
\small
\begin{tabular}{lcccc}
\toprule
\textbf{Model} & \textbf{Biased} & \textbf{Corrected} & \textbf{Biased $p$} & \textbf{Corrected $p$} \\
\midrule
R1-32B & L10 (16\%) & L63 (98\%) & 0.050 (NS) & ${<}0.001$ \\
QwQ-32B & L2 (3\%) & L60 (94\%) & 0.346 (NS) & ${<}0.001$ \\
\bottomrule
\end{tabular}
\caption{Precision significance at biased vs.\ corrected layers. Layer selection metric choice entirely determines whether the HS advantage is detected.}
\label{tab:biased_vs_corrected}
\end{table}

LOO-CV confirms selection stability: the dominant layer is chosen in 65--68\% of folds regardless of the metric used, ruling out validation set overfitting.

\paragraph{Anti-circularity analysis.}
A potential concern is that selecting layers by threshold-dependent metrics (precision, FPR) and then evaluating on the same metrics introduces circularity.
We address this by computing HS vs.\ text comparisons at BA-optimal layers (a threshold-free criterion): all four models pass Holm--Bonferroni on all four evaluation metrics (16/16), with precision advantages of $+$13.7 to $+$25.4$\pp$.
The HS advantage is a genuine modality property, not a layer selection artifact.


\subsection{HS vs.\ Text at Corrected Layers}
\label{sec:hs_vs_text}

Table~\ref{tab:hb_summary} presents the complete Holm--Bonferroni analysis.
All 12 tests pass step-down correction at $\alpha = 0.05$, up from 5/12 with biased layers.

\begin{table}[t]
\centering
\small
\begin{tabular}{rlrrl}
\toprule
\textbf{Rank} & \textbf{Model} & \textbf{Metric} & \textbf{$\Delta$} & \textbf{$p$} \\
\midrule
1 & R1-8B & Prec & $+$22.4$\pp$ & ${<}0.001$ \\
2 & OT-7B & BA & $+$10.7$\pp$ & ${<}0.001$ \\
3 & OT-7B & Prec & $+$21.9$\pp$ & ${<}0.001$ \\
4 & OT-7B & FPR & $-$17.1$\pp$ & ${<}0.001$ \\
5 & QwQ & BA & $+$15.4$\pp$ & ${<}0.001$ \\
6 & QwQ & Prec & $+$23.6$\pp$ & ${<}0.001$ \\
7 & R1-32B & BA & $+$18.7$\pp$ & ${<}0.001$ \\
8 & R1-32B & Prec & $+$15.4$\pp$ & ${<}0.001$ \\
9 & QwQ & FPR & $-$9.8$\pp$ & 0.004 \\
10 & R1-8B & BA & $+$8.2$\pp$ & 0.005 \\
11 & R1-32B & FPR & $-$14.6$\pp$ & 0.007 \\
12 & R1-8B & FPR & $-$7.4$\pp$ & 0.044 \\
\bottomrule
\end{tabular}
\caption{Holm--Bonferroni step-down results (12/12 pass, $\alpha=0.05$). $\Delta$ is HS minus text (positive = HS better for BA/Prec, negative = HS better for FPR). All tests use 10{,}000 bootstrap resamples with null-centered two-sided procedure.}
\label{tab:hb_summary}
\end{table}

Precision advantage is remarkably consistent across models: $+$15.4 to $+$23.6$\pp$.

\paragraph{5-seed robustness.}
Across five random seeds, OT-7B, QwQ-32B, and R1-32B each pass 5/5 Holm--Bonferroni tests.
R1-8B passes 3/5 (failures are FPR-specific).
Precision delta never reverses across any of the 20 seed--model combinations.

\paragraph{Control task selectivity.}
Selectivity (real minus random-label probe accuracy) ranges from 7.5$\pp$ (R1-32B) to 18.3$\pp$ (OT-7B), confirming probes exploit safety-specific signal rather than positional or structural artifacts.

\paragraph{Position mediation.}
Relative position ($t/T$) is a significant component of probe performance.
After residualizing position from features, 3/4 models retain significant HS advantage (OT-7B $+$6.9$\pp$, R1-8B $+$6.4$\pp$, QwQ $+$3.7$\pp$; all $p < 0.011$).
R1-32B's advantage is almost entirely position-mediated (residualized gap $+$1.9$\pp$, $p = 0.914$), consistent with its other caveats (W=1 only significant, 25\% coverage).

\paragraph{32B model-specific caveats.}
The two 32B models exhibit divergent practical profiles despite both showing significant HS advantage:

\begin{itemize}[nosep,leftmargin=*]
\item \textbf{R1-32B}: Significant at W=1 only (3/3 Holm--Bonferroni). Excellent lead time ($+$11.7 to $+$13.3 steps) but narrow coverage (25\%). Dimensionality (51{,}200$\times W$ features for $n=20$) causes overfitting at $W \geq 3$.
\item \textbf{QwQ-32B}: Significant across all 15 window sizes tested. Robust coverage (74\%) but HS detects \emph{later} than text ($-$1.2 to $+$1.9 steps vs.\ text $+$3.7--12.3).
\end{itemize}


\subsection{Temporal Context Improvement}
\label{sec:temporal}

Temporal context universally improves text probes across all four models ($+$3--12$\pp$ BA).
For HS probes, 3/4 models benefit from larger temporal windows; OT-7B shows early saturation with BA peaking at $W=3$ (0.688) while precision and FPR continue improving.

Table~\ref{tab:temporal_32b} reports temporal evaluation on the 32B models at corrected layers.
R1-32B achieves its best BA (0.737) at $W=10$ with $+$1.9$\pp$ over static, but only $W=1$ passes Holm--Bonferroni---a power limitation rather than temporal ineffectiveness ($n=20$ with 51{,}200-dimensional features).
QwQ-32B peaks at $W=5$ (BA 0.761, $+$3.1$\pp$ over static) with all windows $W=1$--15 passing significance tests.

\begin{table}[t]
\centering
\small
\begin{tabular}{llccccc}
\toprule
\textbf{Model} & \textbf{W} & \textbf{BA} & \textbf{Prec} & \textbf{FPR} & \textbf{Cov} & \textbf{HB} \\
\midrule
R1-32B & 1 & .826 & .603 & .235 & 25\% & 3/3 \\
R1-32B & 3 & .738 & .560 & .306 & 80\% & 0/3 \\
R1-32B & 10 & .737 & .609 & .234 & 70\% & 0/3 \\
QwQ & 1 & .730 & .544 & .228 & 79\% & 3/3 \\
QwQ & 3 & .745 & .565 & .214 & 74\% & 3/3 \\
QwQ & 5 & .761 & .573 & .217 & 74\% & 3/3 \\
\bottomrule
\end{tabular}
\caption{Temporal probe evaluation on 32B models at corrected layers. HB = Holm--Bonferroni tests passed (out of 3 metrics).}
\label{tab:temporal_32b}
\end{table}

\paragraph{Window size ablation.}
Across all four models, coverage monotonically decreases with $W$---larger windows create more selective features, reducing both false and true positives.
This is a precision--coverage tradeoff, not evidence against temporal context.
R1-8B precision increases from 0.693 ($W=1$) to 0.822 ($W=25$) while coverage drops from 72.1\% to 48.8\%.

\begin{table}[t]
\centering
\small
\begin{tabular}{rcccc}
\toprule
\textbf{W} & \textbf{R1-8B} & \textbf{OT-7B} & \textbf{QwQ} & \textbf{R1-32B} \\
\midrule
1 & .778 & .672 & .730 & .718 \\
3 & .790 & .688 & .745 & .725 \\
5 & .781 & .675 & .761 & .731 \\
10 & .793 & .660 & .747 & .737 \\
15 & .802 & .650 & .741 & .721 \\
25 & .810 & .635 & .762 & .703 \\
\bottomrule
\end{tabular}
\caption{Balanced accuracy by window size and model (HS probes at corrected layers). OT-7B peaks at $W=3$; other models benefit from larger windows.}
\label{tab:window_ablation}
\end{table}


\subsection{Detection Timing Profiles}
\label{sec:timing}

Lead time---the number of steps before the commitment point at which the probe first fires---varies substantially across models (Table~\ref{tab:lead_time}).

\begin{table}[t]
\centering
\small
\begin{tabular}{lcccc}
\toprule
\textbf{Model} & \textbf{HS LR} & \textbf{Text} & \textbf{GRU} & \textbf{Note} \\
\midrule
R1-8B & $\sim$0 & $+$10.3* & $+$4.5$\pm$3.0 & HS at CP \\
OT-7B & $+$21.6$\dagger$ & $+$45.5$\dagger$ & $+$33.6$\pm$4.8 & $\dagger$length \\
QwQ & $-$1 to $+$2 & $+$4--12 & $+$15.9 & HS late \\
R1-32B & $+$12--13 & $+$4--9 & $+$8.8 & HS early \\
\bottomrule
\end{tabular}
\caption{Lead time profiles (steps before CP). *Text lead inflated by 20.5\% FPR. $\dagger$OT-7B absolute values inflated by long traces ($\sim$80 sentences).}
\label{tab:lead_time}
\end{table}

On R1-8B, the HS probe detects \emph{at} the commitment point (lead $\approx 0$); only the GRU achieves modest positive lead ($+$4.5$\pm$3.0, worst seed $+$0.9).
Text's earlier detection ($+$10.3) is an artifact of its higher FPR (20.5\%)---it is trigger-happy rather than prescient.
R1-32B and QwQ-32B diverge: R1-32B shows excellent HS lead ($+$12--13 steps) while QwQ's HS detects \emph{later} than text, illustrating that the optimal detection modality depends on the specific deployment requirement (early warning vs.\ broad accuracy).

\paragraph{Oracle decomposition.}
An oracle detector using ground-truth label sequences achieves $+$31 steps lead with 100\% detection (on R1-8B).
The gap arises because the oracle exploits clean consecutive-label patterns, whereas probes must infer labels from noisy hidden-state features where confidence fluctuates around the CP.
This gap quantifies a fundamental limitation: hidden states encode \emph{current} safety state but poorly predict \emph{future} trajectory.

\subsection{Comprehensive Negative Results}
\label{sec:negative}

Six alternative detection approaches fail across all four models, establishing important negative baselines.

\paragraph{Mahalanobis distance.}
TrajGuard-style detection using Mahalanobis distance on PCA-reduced hidden states achieves near-zero separation between safe and post-commitment distributions across all models (Table~\ref{tab:mahalanobis}).
The failure mode is fundamental: safety commitment signal resides in minor PCA components, not principal directions.

\begin{table}[t]
\centering
\small
\begin{tabular}{lccccl}
\toprule
\textbf{Model} & \textbf{Layer} & \textbf{PCA} & \textbf{Safe} & \textbf{Post} & \textbf{Sep} \\
\midrule
R1-8B & L14 & 64 & 6.66 & 6.73 & 0.07 \\
OT-7B & L16 & 64 & 7.68 & 7.71 & 0.03 \\
QwQ & L60 & 16 & 3.22 & 3.42 & 0.20 \\
R1-32B & L63 & 64 & 6.73 & 6.88 & 0.15 \\
\bottomrule
\end{tabular}
\caption{Mahalanobis distance discrimination. Separation $\approx 0$ across all models; false alarm rate $\geq$ true positive rate at every threshold.}
\label{tab:mahalanobis}
\end{table}

\paragraph{BOCPD and CUSUM.}
Bayesian Online Change-Point Detection \cite{adams2007bocpd} achieves 0\% detection on 3/4 models (OT-7B, QwQ, R1-32B); only R1-8B reaches 93\% detection but with lower lead time ($+$3.5) than simple thresholding ($+$4.9).
The Gaussian changepoint assumption is violated by the abrupt commitment transitions in most models.
CUSUM performs universally worse than simple thresholds.

\begin{table}[t]
\centering
\small
\begin{tabular}{lcccl}
\toprule
\textbf{Model} & \textbf{BOCPD} & \textbf{Lead} & \textbf{FCP Lead} & \textbf{Winner} \\
\midrule
R1-8B & 93\% & $+$3.5 & $+$4.9 & FCP \\
OT-7B & 0\% & --- & $+$7.5 & FCP \\
QwQ & 0\% & --- & $+$9.4 & FCP \\
R1-32B & 0\% & --- & $+$2.6 & FCP \\
\bottomrule
\end{tabular}
\caption{BOCPD detection rate and lead time vs.\ fixed-crossing-point (FCP) baseline. BOCPD fails on 3/4 models due to violated Gaussian assumptions.}
\label{tab:bocpd}
\end{table}

\paragraph{Cross-model transfer.}
Zero-shot cross-model transfer (using PCA projection for dimension alignment) yields near-chance accuracy across all pairs (Table~\ref{tab:transfer}).
The best transfer (QwQ $\to$ R1-32B: 0.673) falls well below the weakest in-domain result (0.652).
Safety commitment representations are model-specific, ruling out universal safety probes.

\begin{table}[t]
\centering
\small
\begin{tabular}{lcccc}
\toprule
\textbf{Train $\to$ Test} & \textbf{R1-8B} & \textbf{R1-32B} & \textbf{QwQ} \\
\midrule
R1-8B & 0.823 & 0.522 & 0.595 \\
R1-32B & 0.583 & 0.652 & 0.550 \\
QwQ & --- & 0.673 & 0.717 \\
\midrule
In-domain & 0.823 & 0.652 & 0.717 \\
\bottomrule
\end{tabular}
\caption{Cross-model transfer accuracy. All off-diagonal entries are at or near chance, indicating model-specific safety representations.}
\label{tab:transfer}
\end{table}

\paragraph{GRU vs.\ linear probes.}
On 32B models, GRU probes do \emph{not} outperform temporal logistic regression: QwQ GRU BA 0.724 vs.\ LR 0.761 ($-$3.7$\pp$); R1-32B GRU 0.675 vs.\ LR 0.737 ($-$6.2$\pp$).
On 7--8B models, GRU provides modest improvement (R1-8B $+$0.7$\pp$, OT-7B $+$5.8$\pp$) with longer lead times at the cost of higher FPR.
Linear probes are sufficient for safety commitment detection.

\begin{table}[t]
\centering
\small
\begin{tabular}{lccccc}
\toprule
\textbf{Model} & \textbf{LR BA} & \textbf{GRU BA} & \textbf{$\Delta$} & \textbf{GRU Lead} & \textbf{LR Lead} \\
\midrule
R1-8B & .802 & .809 & $+$0.7 & $+$4.5 & $-$0.1 \\
OT-7B & .688 & .746 & $+$5.8 & $+$33.6 & $+$21.6 \\
QwQ & .761 & .724 & $-$3.7 & $+$15.9 & $\sim$0 \\
R1-32B & .737 & .675 & $-$6.2 & $+$9.1 & $+$12.9 \\
\bottomrule
\end{tabular}
\caption{GRU vs.\ temporal LR comparison. GRU underperforms LR on 32B models; modest advantage on 7--8B trades higher FPR for longer lead.}
\label{tab:gru_vs_lr}
\end{table}


\subsection{HS vs.\ Text Baseline Detail}
\label{sec:text_baseline}

Table~\ref{tab:hs_text_detail} provides the full HS vs.\ text comparison on R1-8B.
HS probes achieve $+$15--18$\pp$ accuracy over text, $+$10$\pp$ precision, and $-$19.4$\pp$ FPR.
Text achieves non-trivial performance (67.8\% BA)---CoT text does carry safety information---but HS provides substantial unique value from sub-textual signal.

\begin{table}[t]
\centering
\small
\begin{tabular}{lccccc}
\toprule
\textbf{Method} & \textbf{BA} & \textbf{Prec} & \textbf{FPR$_\text{pre}$} & \textbf{Lead} & \textbf{Cov} \\
\midrule
HS Static & .768 & .622 & .162 & $+$3.0 & .605 \\
HS Temp W15 & .726 & .583 & .160 & $+$4.9 & .651 \\
Text Static & .661 & .448 & .356 & $+$4.3 & .605 \\
Text Temp W15 & .678 & .487 & .272 & $+$11.5 & .651 \\
GRU & .816 & .648 & .188 & $+$5.3 & .744 \\
\bottomrule
\end{tabular}
\caption{Full method comparison on R1-8B. HS probes outperform text on BA, precision, and FPR. Text's higher lead time is inflated by higher FPR.}
\label{tab:hs_text_detail}
\end{table}


\subsection{OpenThinker-7B: Cleanest Single-Model Evidence}
\label{sec:ot7b}

OT-7B (Qwen-7B, STILL-trained, $n=228$ traces) provides the cleanest evidence for HS advantage: 3/3 Holm--Bonferroni pass with the largest effect sizes, highest selectivity (18.3$\pp$), and 5/5 seed robustness.
At layer 16 with $W=3$: BA 0.688 (vs.\ text 0.581), precision 0.466 (vs.\ 0.254), FPR 10.0\% (vs.\ 27.1\%).
Its temporal saturation at $W=3$ is a genuine model characteristic confirmed on multiple layers, with precision and FPR continuing to improve at larger $W$ even as BA declines.

%----------------------------------------------------------------------
\section{Discussion}
\label{sec:discussion}
%----------------------------------------------------------------------

\paragraph{Layer selection as researcher degree of freedom.}
Our primary finding---that the metric \texttt{bal\_acc $\times$ crossing\_rate} systematically biases layer selection toward shallow layers---has implications beyond this study.
Any investigation using threshold-frequency--contaminated metrics to select layers or hyperparameters risks the same systematic bias.
The 5-metric $\times$ 4-model ablation constitutes a reproducible audit protocol that we recommend as standard practice.

\paragraph{Oracle gap and fundamental limits.}
The $+$31-step oracle lead vs.\ $+$4.5-step GRU lead reveals that hidden states encode current safety state (80\% BA) but poorly predict future trajectory.
Pre-CP safety streaks average 14.6 consecutive safe labels; probes cannot distinguish ``safely exploring'' from ``about to commit'' based on snapshot features alone.
Closing this gap likely requires predictive probes trained on future safety state or attention pattern analysis.

\paragraph{Safe-trace false positive rates.}
All methods exhibit high false positive rates on safe traces (11.1\%--62\% across models and methods).
OT-7B is the exception at 11.1\% (HS static); most model--method combinations range from 39\% to 62\%.
These results establish baselines for future improvement rather than production-ready thresholds.
HS consistently achieves lower safe FPR than text in 3/4 models.

\begin{table}[t]
\centering
\small
\begin{tabular}{lcccc}
\toprule
\textbf{Model} & \textbf{HS Stat} & \textbf{HS Temp} & \textbf{Text Stat} & \textbf{Text Temp} \\
\midrule
R1-8B & 39.0\% & 39.2\% & 50.2\% & 38.5\% \\
OT-7B & 11.1\% & 17.5\% & 30.2\% & 46.0\% \\
QwQ & 50.0\% & 45.0\% & 32.5\% & 42.5\% \\
R1-32B & 23.9\% & 23.9\% & 31.0\% & 31.0\% \\
\bottomrule
\end{tabular}
\caption{False positive rate on safe traces. Most methods exhibit 39--62\% safe FPR; OT-7B HS static is the exception at 11.1\%.}
\label{tab:safe_fpr}
\end{table}

\paragraph{32B divergence.}
R1-32B and QwQ-32B exhibit opposing practical profiles: R1-32B offers excellent lead time ($+$13 steps) with narrow coverage (25\%), while QwQ offers robust significance with poor lead time.
This divergence is not a weakness but a genuine finding: safety commitment dynamics at 32B scale are model-specific, and deployment guidance must be per-model.

%----------------------------------------------------------------------
\section{Limitations}
\label{sec:limitations}
%----------------------------------------------------------------------

Our study has three principal limitations.
First, sample sizes for 32B models are small ($n=19$--20 test traces), limiting statistical power---R1-32B achieves significance only at $W=1$, and the residualized position analysis is underpowered (MDE 11.2$\pp$ vs.\ observed gap 1.9$\pp$).
Second, safe-trace false positive rates (39--62\% for most configurations) are too high for production deployment; our results establish baselines rather than deployable thresholds.
Third, GRU probes on 32B models were evaluated at biased layers and should be re-evaluated at corrected layers in future work.

%----------------------------------------------------------------------
\section{Conclusion}
\label{sec:conclusion}
%----------------------------------------------------------------------

We present the first systematic evaluation of safety commitment detection in reasoning models, spanning four models, two architectures, and two scales.
Our primary methodological contribution is the discovery that the layer selection metric \texttt{bal\_acc $\times$ crossing\_rate} systematically biases toward shallow layers, entirely masking hidden-state probe advantages at 32B scale and flipping results from 5/12 to 12/12 significant under Holm--Bonferroni correction.
At correctly selected layers, hidden-state probes significantly outperform text-based probes across all four models (precision $+$15--24$\pp$), with temporal context providing universal improvement.
The comprehensive failure of six alternative methods---Mahalanobis distance, BOCPD, CUSUM, cross-model transfer, CCS, and RepE---establishes that supervised temporal probing with principled layer selection is the most effective current approach.
These findings provide both a methodological caution (layer selection metrics matter) and an empirical foundation for reasoning-model safety monitoring.

%----------------------------------------------------------------------
% References
%----------------------------------------------------------------------

\begin{thebibliography}{10}

\bibitem{adams2007bocpd}
R.~P. Adams and D.~J.~C. MacKay.
Bayesian online changepoint detection.
{\em arXiv:0710.3742}, 2007.

\bibitem{burns2023ccs}
C.~Burns, H.~Ye, D.~Klein, and J.~Steinhardt.
Discovering latent knowledge in language models without supervision.
In {\em ICLR}, 2023.

\bibitem{huang2025selfJailbreaking}
Y.~Huang et~al.
Self-jailbreaking in reasoning models.
{\em arXiv preprint}, 2025.

\bibitem{li2024probing}
K.~Li, O.~Patel, F.~Vi\'{e}gas, H.~Pfister, and M.~Wattenberg.
Inference-time intervention: Eliciting truthful answers from a language model.
In {\em NeurIPS}, 2024.

\bibitem{peng2025siren}
Y.~Peng et~al.
{SIREN}: Internal representation monitoring for jailbreak detection.
{\em arXiv preprint}, 2025.

\bibitem{wang2025safeRemind}
X.~Wang et~al.
{SafeRemind}: Chain-of-thought safety monitoring.
{\em arXiv preprint}, 2025.

\bibitem{wen2025harmThoughts}
Y.~Wen et~al.
{HarmThoughts}: Harmful thoughts in reasoning model chains of thought.
{\em arXiv preprint}, 2025.

\bibitem{zhang2025trajguard}
Z.~Zhang et~al.
{TrajGuard}: Trajectory-level jailbreak detection via hidden state analysis.
{\em arXiv preprint}, 2025.

\bibitem{zheng2025safeThink}
C.~Zheng et~al.
{SafeThink}: Safety monitoring for chain-of-thought reasoning.
{\em arXiv preprint}, 2025.

\bibitem{zou2023repE}
A.~Zou, L.~Phan, et~al.
Representation engineering: A top-down approach to {AI} transparency.
{\em arXiv:2310.01405}, 2023.

\end{thebibliography}

%----------------------------------------------------------------------
\appendix
\section{16-Label Taxonomy Supervision}
\label{app:taxonomy}
%----------------------------------------------------------------------

We test whether fine-grained safety labels (16 categories) improve over binary supervision.
On 7--8B models, 16-label training yields modest coverage improvement ($+$5$\pp$) with marginal precision gains ($+$1--3$\pp$).
32B models are inconclusive due to sample size limitations (16 classes dilute per-class counts).

\begin{table}[h]
\centering
\small
\begin{tabular}{llcccc}
\toprule
\textbf{Model} & \textbf{Supervision} & \textbf{BA} & \textbf{Prec} & \textbf{Cov} & \textbf{Lead} \\
\midrule
R1-8B & binary & .802 & .777 & .512 & $-$0.1 \\
R1-8B & 16-label & .820 & .795 & .558 & $+$0.8 \\
OT-7B & binary & .688 & .466 & .532 & 21.6 \\
OT-7B & 16-label & .700 & .485 & .580 & 22.3 \\
\bottomrule
\end{tabular}
\caption{Binary vs.\ 16-label supervision on 7--8B models. Improvement is modest; temporal windowing is a more effective lever.}
\label{tab:taxonomy}
\end{table}

%----------------------------------------------------------------------
\section{Ensemble CV Layer Selection}
\label{app:ensemble_cv}
%----------------------------------------------------------------------

We validate multi-metric ensemble layer selection via 5-fold cross-validation on 7--8B models.
HS outperforms text in 10/10 folds with mean precision delta $+$17.4$\pm$2.6$\pp$ (R1-8B) and $+$16.4$\pm$3.3$\pp$ (OT-7B).
All fold-level comparisons are significant ($p_{\max}=0.002$).

\begin{table}[h]
\centering
\small
\begin{tabular}{llcccccl}
\toprule
\textbf{Model} & \textbf{Metric} & \textbf{F1} & \textbf{F2} & \textbf{F3} & \textbf{F4} & \textbf{F5} & \textbf{Mean$\pm$Std} \\
\midrule
R1-8B & $\Delta$Prec & +.19 & +.16 & +.18 & +.15 & +.19 & $.174\pm.026$ \\
OT-7B & $\Delta$Prec & +.20 & +.13 & +.17 & +.16 & +.16 & $.164\pm.033$ \\
\bottomrule
\end{tabular}
\caption{Ensemble CV precision delta (HS$-$Text) per fold. HS advantage is consistent across all folds.}
\label{tab:ensemble_cv}
\end{table}

\end{document}
